% ===========================================================================
% CHAPTER 6
% ===========================================================================

\chapter{Certificate Transparency Renewal and Discovery Pipeline}
\label{chap:ct-renewal-pipeline}

Certificate Authorities (CAs) participating in the Web Public Key
Infrastructure are required by the CA/Browser Forum Baseline Requirements to
submit every issued certificate to multiple public Certificate Transparency
(CT) logs~\cite{rfc9162}. These logs maintain append-only, cryptographically
verifiable records of all certificate issuance events. The public visibility
created by this obligation makes CT logs a valuable observational source for
monitoring certificate activity, including detecting newly issued certificates,
identifying renewals of existing deployments, and observing changes in the
certificate ecosystem over time. The pipeline described in this chapter
converts these raw CT observations into a continuously maintained dataset by
collecting domain events in real time, periodically separating renewed
certificates from those belonging to previously unseen domains, acquiring
current endpoint data for both groups, and incorporating the results into the
analytical dataset.

\section{Continuous CT Observation}
\label{sec:ct-collection-configuration}

The Certificate Transparency observation subsystem supplies the domain-level
events that drive the maintenance workflow. It comprises three cooperating
artifacts: a CT log listener, a consumer process, and a per-log position
index.

The CT log listener maintains persistent connections to a configured set of
public logs---operated by Google, Cloudflare, Let's Encrypt, DigiCert,
Sectigo, and others. As certificates are submitted to these logs, the
listener decodes each issuance event and exposes the observed domain sets
through a local WebSocket endpoint, delivering a stream of domain and Subject
Alternative Name (SAN) observations extracted from each processed
certificate.

A consumer process maintains a persistent WebSocket connection to this
endpoint. For each received message, it strips wildcard markers,
deduplicates the domain set within that message, and accumulates unique
observations in an in-memory buffer. Each buffered entry records the list of
observed domains together with a processing flag initially set to
\texttt{false}. When the buffer reaches a configured threshold, the
accumulated records are flushed to MongoDB in a single batch write. A unique
index on the \texttt{domains} field ensures that repeat observations of the
same domain set are silently rejected at the database level without
disrupting the ingestion flow.

Two complementary persistence mechanisms support uninterrupted operation
across process restarts. The per-log position index records the most recently
processed entry position for each monitored CT log, allowing the listener to
resume observation from that position rather than reprocessing the complete
log history. Separately, the consumer records ingestion metadata---batch
counts, processing timestamps, and run status---in a dedicated MongoDB
collection. The processing flag on each buffered domain set serves as an
additional tracking mechanism: it is set to \texttt{true} when the
corresponding set has been considered by the renewal and discovery workflow,
preventing the same observation from triggering repeated processing
cycles.

This design supports continuous, low-latency observation of CT log activity.
On graceful shutdown, any remaining entries in the in-memory buffer are
flushed to MongoDB before the process terminates. The buffer is not a durable
message queue and does not guarantee exactly-once delivery for every domain
observation, but this is acceptable given that the subsequent maintenance
cycle processes the collected data in bulk and can tolerate occasional
duplicate or missed observations within practical bounds.

\section{The Periodic Maintenance Cycle}
\label{sec:eight-stage-orchestrator}

While the observation subsystem runs continuously, the maintenance operations
that convert collected observations into dataset updates execute as a
periodic, orchestrator-driven cycle. An orchestration component sequences
several discrete stages in a defined order, ensuring that each completed
cycle produces a coherent, self-consistent update to the certificate dataset
and the corresponding analytical results. Figure~\ref{fig:ct-renewal-discovery-branches}
illustrates the overall flow.

\begin{figure}[H]
  \centering
  \begin{tikzpicture}[
    font=\scriptsize,
    >=Latex,
    box/.style={draw, rounded corners=2pt, align=center, text width=2.8cm, minimum height=0.8cm, fill=black!5},
    store/.style={draw, cylinder, shape border rotate=90, aspect=0.22, align=center, text width=2.8cm, minimum height=0.9cm, fill=black!8},
    main/.style={draw, very thick, cylinder, shape border rotate=90, aspect=0.22, align=center, text width=2.8cm, minimum height=0.95cm, fill=black!10},
    flow/.style={-{Latex[length=2mm]}, line width=0.5pt},
    dashflow/.style={-{Latex[length=2mm]}, line width=0.5pt, dashed}
  ]
    \node[box] (ct-logs) at (0,2.8) {Public CT logs};
    
    \node[box] (listener) at (0,1.6) {CT log listener\\and domain ingestion};
    \node[store] (live-db) at (0,-0.45) {Live domain database\\(continuously updated)};

    \node[box] (renew-select) at (-3.8,-2.6) {\hyphenpenalty=10000\exhyphenpenalty=10000 Match observations with\\known domains};
    \node[box] (new-select) at (3.8,-2.6) {\hyphenpenalty=10000\exhyphenpenalty=10000 Select previously unseen\\candidate domains};

    \node[box] (renew-crawl) at (-3.8,-4.65) {Crawl renewal\\candidates};
    \node[box] (new-crawl) at (3.8,-4.65) {Crawl discovery\\candidates};

    \node[store] (renew-db) at (-3.8,-6.45) {Renewal staging\\certificates};
    \node[store] (new-db) at (3.8,-6.45) {Discovery staging\\certificates};

    \node[box] (renew-merge) at (-3.8,-8.0) {Checkpointed replacement\\by domain};
    \node[box] (new-merge) at (3.8,-8.0) {Insert accepted records\\and extend domain list};
    \node[main] (main-db) at (0,-9.6) {Updated certificate\\dataset};

    \node[box] (dashboard) at (0,-11.0) {Interactive dashboard};

    \draw[flow] (ct-logs) -- (listener);
    \draw[flow] (listener) -- (live-db);
    \draw[dashflow] (live-db.south) -- (renew-select.north);
    \draw[dashflow] (live-db.south) -- (new-select.north);
    \draw[flow] (renew-select) -- (renew-crawl);
    \draw[flow] (renew-crawl) -- (renew-db);
    \draw[flow] (renew-db) -- (renew-merge);
    \draw[flow] (renew-merge) -- (main-db);
    \draw[flow] (new-select) -- (new-crawl);
    \draw[flow] (new-crawl) -- (new-db);
    \draw[flow] (new-db) -- (new-merge);
    \draw[flow] (new-merge) -- (main-db);
    \draw[flow] (main-db) -- (dashboard);
  \end{tikzpicture}
  \caption{The two-branch pipeline: CT observations are split into renewal candidates (left, replacing existing records) and discovery candidates (right, inserted as new records), both converging on the updated dataset that feeds the dashboard.}
  \label{fig:ct-renewal-discovery-branches}
\end{figure}

The cycle begins by temporarily stopping the CT observation subsystem so that
the accumulated domain data in the live database remains stable throughout
the maintenance run. The orchestrator then verifies that the global domain
dataset exists and, if necessary, generates it from the primary MongoDB
collection. The core processing stage separates the buffered CT observations
into two categories---domains already present in the maintained dataset
(renewal candidates) and domains that have not been observed before
(discovery candidates)---and processes each category independently through
the certificate acquisition workflow described in
Chapter~\ref{chap:certificate-acquisition}. After acquisition completes, the
newly obtained certificates are merged into the main dataset using different
strategies for each category: renewal results replace existing records,
while discovery results are inserted as new entries. The analytical results
are then recomputed, temporary staging data is removed, and the CT
observation subsystem is restarted to resume collection for the next cycle.

\section{Renewal Detection and New Domain Discovery}
\label{sec:renewal-and-discovery}

Renewal detection determines which of the observed CT domain sets correspond
to certificates already represented in the maintained dataset. For each
domain in the global domain list, the system generates base-domain and
\texttt{www} variations and queries the buffered CT observations in the live
domain database. A unique index on the stored domain sets allows each
observation to be located efficiently by any domain that appears in its SAN
list.

When a matching observation is found, only one representative domain is
selected for the renewal output, even when multiple domains in the dataset
share the same certificate (indicated by identical SAN sets). This prevents
redundant crawling of the same certificate endpoint from multiple domain
entries. The corresponding CT observation is marked as processed so that
subsequent maintenance cycles do not re-examine it. The resulting renewal
list therefore contains one entry per renewed certificate rather than one
entry per matching domain, reducing the acquisition workload while
maintaining full coverage of the observed renewal activity.

The discovery path extends the dataset with certificates for domains that
have not been previously observed. It begins by extracting a bounded number
of unprocessed domain records from the live domain database, writing them to
a temporary input file. The source collection is then dropped, ensuring that
the next observation cycle starts with a clean buffer.

The domain-based crawler is launched against this extracted input and runs to
completion. Each successfully acquired certificate is stored in a staging
database. In the subsequent merge step, the newly acquired certificates are
inserted into the main dataset as new records keyed by domain. Because
modern certificates commonly cover a large number of Subject Alternative
Names (SANs), a single certificate may be acquired independently for several
different domains during a single discovery cycle. No fingerprint-based
deduplication is applied---each domain receives its own independent record
in the dataset, which is the desired behaviour given that the same
certificate may serve multiple distinct web properties. The system tracks
exactly which domains were successfully inserted and appends only those
confirmed domains to the global domain list, maintaining consistency between
the MongoDB records and the domain list used for subsequent acquisition
cycles.

\section{Dataset Merge and Analytical Refresh}
\label{sec:ct-precompute-integration}

The renewal merge transfers the certificates acquired during the renewal
crawl from the staging database into the main certificate dataset. Each
renewal document is written using a replacement operation keyed on the
domain field: if a record for that domain already exists in the main
collection, it is replaced in its entirety; if no record exists, a new
document is inserted. This approach avoids the two-step delete-and-insert
pattern, which would create a window during which the dashboard could read an
incomplete certificate record.

Progress is tracked through a checkpoint stored in a metadata collection
within the main database. The checkpoint records the identifier of the last
successfully processed renewal document and is advanced only after each batch
write completes. If the merge process is interrupted---by a crash, a timeout,
or an explicit stop request---a subsequent run resumes from the stored
checkpoint rather than reprocessing the complete renewal dataset. Because
each replacement operation is idempotent, rewriting the same document
multiple times produces identical results, making partial progress always
safe to resume.

After the certificate dataset has been updated, the analytics process
verifies the required database indexes and recomputes the six materialized
analytical families. Detailed certificate queries and selected time-series
views read directly from the main dataset and therefore reflect the current
state immediately. Materialized dashboard views, by contrast, serve the
latest completed analytical snapshot. The resulting availability is near real
time in the sense that it follows the completed maintenance cycle; it is not
an instantaneous streaming update driven by individual CT events.

\section{Component Map}
\label{sec:ct-component-map}

Table~\ref{tab:ct-components} maps each functional role described in this
chapter to the concrete artifact that implements it. Most artifacts reside in
the \texttt{ct-logs-renewal-pipeline/} directory; the orchestrator additionally
invokes the domain-based crawler in \texttt{ssl-certificates-crawler/} and the
analytical refresh script \texttt{run-generic.py} in
\texttt{dashboard/backend/pre-compute-scripts/}. All stages are coordinated by
the orchestrator described in Section~\ref{sec:eight-stage-orchestrator}.

\renewcommand{\arraystretch}{1.3}

\begin{table}[H]
\centering
\small
\begin{tabularx}{\textwidth}{@{}p{3.2cm} p{4.3cm} X@{}}
\toprule
\textbf{Functional role} & \textbf{Artifact} & \textbf{Responsibility} \\
\midrule

CT log listener &
\texttt{certstream-server-go} &
Maintains persistent connections to public CT logs and exposes domain events over a local WebSocket \\

Observation consumer &
\texttt{go-server.py} &
Subscribes to the listener, buffers observations in memory, and batch-writes unique domain sets to MongoDB \\

Log position index &
\texttt{ct\_index.json} &
Records the last-seen position per CT log for resumption after a restart \\

Orchestrator &
\texttt{main.py} &
Sequences the stages of each maintenance cycle \\

Master domain list &
\texttt{global-dataset.csv} &
Authoritative list of known domains, generated by \texttt{fetch-domains-names.py} \\

Renewal detection &
\texttt{data-renew.py} &
Matches buffered observations against known domains \\

Renewal merge &
\texttt{data-renew-merge.py} &
Performs checkpointed, domain-keyed replacement of renewal certificates into the main dataset \\

Discovery acquisition &
\texttt{new-data.py} &
Crawls previously unseen domains and inserts new certificates into the main dataset\\

Analytical refresh &
\texttt{run-generic.py} &
Recomputes the required indexes and the six materialized analytical families \\

\bottomrule
\end{tabularx}
\caption{Mapping of pipeline roles to implementing artifacts.}
\label{tab:ct-components}
\end{table}

\renewcommand{\arraystretch}{1}


\section{Reliability and Correctness Considerations}
\label{sec:ct-failure-limitations}

The pipeline incorporates several stage-level recovery mechanisms. CT-log
indices in the observation service allow collection to resume from a known
position. Persistent processing flags prevent the same domain set from
triggering repeated processing cycles. Acquisition task states, batch
checkpoints, process-state tracking, and failure logs allow individual stages
to be resumed or repeated after interruption without restarting the complete
workflow from the beginning. In practice the orchestrator is launched on a
fixed schedule through the operating system's task scheduler (cron on
Linux/macOS, Task Scheduler on Windows), with each invocation running its
stages to completion before the next cycle begins.

The complete maintenance cycle is not, however, enclosed within a single
transaction. CT observation, temporary staging, domain-list updates,
certificate merges, analytical recomputation, and subsystem restarts each
remain separate operations. The renewal checkpoint must remain consistent
with the lifetime of its staging data---if the staging database is cleared
before the checkpoint is advanced, the merge cannot be resumed. Concurrent
maintenance cycles also require explicit operational coordination, because
overlapping runs would read from and write to the same staging and main
datasets without isolation guarantees.

Consequently, the platform provides a practical, automated dataset maintenance
process that significantly reduces the manual effort required to keep the
certificate collection current. It does not claim global exactly-once
processing for individual CT events, instantaneous end-to-end updates in
response to every observed issuance, or experimentally measured freshness
bounds. These characteristics are appropriate for the intended research and
analytical use case, where periodic, coherent snapshots of the certificate
ecosystem support meaningful analysis without requiring strict real-time
event processing.
